\pdfminorversion=4
\documentclass[]{article}
\usepackage[utf8]{inputenc}
\usepackage{amssymb,latexsym,amsmath}
\usepackage[a4paper,top=3cm,bottom=2cm,left=3cm,right=3cm,marginparwidth=1.75cm]{geometry}
\usepackage{graphicx}
\usepackage[colorlinks=true, allcolors=blue]{hyperref}
\begin{document}
 
\section{Daten vom Industrie Controller mit der VR Brille synkronisieren}
 
\subsection{OPC UA Client test in Standalone APK}
 
\textbf{Description:} \\
In einem  Unity projekt (Zu finden im Sandbox Folder) haben wir einen auf PC funktionierenden Build erstellt, mitdem unser Programm mit dem Industrie Controller kommunizieren kann. Das builden einer Standalone apk hat funktioniert, die Kommunikation zwischen Industrie Kontroller und VR Brille allerdings nicht
 
 
\subsection{OPC UA Communication Failure in APK Build}
 
\textbf{Description:} \\
Da OPC UA auf Standalone build nicht funktioniert, haben wir das script umgebaut, sodass Kommunikation jetzt über TCP funktioniert.
 
\subsection{TCP Read-Funktion (Statusabfrage)}
 
\textbf{Description:} \
Die Kommunikation über TCP ermöglicht das Auslesen (Read) des aktuellen Zustands der Lampen vom Industrie-Controller.
 
Nach erfolgreicher Verbindung zum Server wird kontinuierlich auf eingehende Nachrichten gewartet. Diese werden zeilenweise empfangen und in eine Queue geschrieben, um sie im Hauptthread von Unity weiterzuverarbeiten.
Wenn ein Read Gesendet Wird: Server antwortet mit Zustandswerten der Lampen im folgenden Format:
 
\begin{verbatim}
room1/Lampe=True
room2/Lampe=False
room3/Lampe=True
\end{verbatim}


\subsection{TCP Write-Funktion (Befehle senden)}
\textbf{Description:} \
Die Write-Funktion ermöglicht das Senden von Steuerbefehlen über TCP an den Industrie-Controller, um den Zustand einzelner Lampen zu setzen. Befehle werden als formatierte Strings über den bestehenden TCP-Stream gesendet.


\begin{verbatim}
::room1:SwitchValueGL=True
::room2:SwitchValueGL=False
::room3:SwitchValueGL=True
\end{verbatim}
\break
\subsection{Lampen-Synchronisation zwischen Unity und Industrie-Controller}
\textbf{Description:} \
Description:
Die Synchronisation stellt sicher, dass der Lampenzustand in Unity jederzeit mit dem tatsächlichen Zustand am Industrie-Controller übereinstimmt. Dazu wird zyklisch ein Read-Befehl gesendet, die Antwort geparst und der interne Unity-Zustand entsprechend aktualisiert. Ändert der Nutzer in Unity den Zustand einer Lampe, wird ein Write-Befehl gesendet und anschließend der Zustand per Read verifiziert.

\begin{verbatim}
Unity sendet konstanten READ
        ↓
Controller antwortet mit Zustandswerten
        ↓
Unity parst Antwort → aktualisiert lokalen Lampenzustand
        ↓
Nutzer ändert Lampenzustand in Unity
        ↓
Unity sendet WRITE (::room1:SwitchValueGL=True)
        ↓
Unity sendet erneut READ zur Verifikation
        ↓
Lokaler Zustand wird aktualisiert
\end{verbatim}

\end{document}